\documentclass[journal]{IEEEtran}
\pdfminorversion=7

\usepackage{amsmath}
\usepackage{array}
\usepackage{booktabs}
\usepackage{graphicx}
\usepackage{tabularx}
\usepackage{url}
\setlength{\emergencystretch}{2em}
\newcolumntype{Y}{>{\raggedright\arraybackslash}X}

\newcommand{\ownedfigurebox}[3]{%
\includegraphics[width=0.92\linewidth]{#2}}

\begin{document}

\title{Protocol-Aware Validation of COVID-19 Respiratory-Audio Screening Under Temporal Drift, External Dataset Shift, and Metadata Confounding}

\author{Author~Names~Withheld%
\thanks{Author affiliations are withheld for review.}%
}

\markboth{IEEE Journal of Biomedical and Health Informatics}%
{Validation Under Shift for COVID-19 Respiratory Audio}

\maketitle

\begin{abstract}
Respiratory-audio screening for COVID-19 is attractive because cough, breathing, and speech can be collected remotely with commodity devices. However, an internal classification score is not sufficient evidence of biomedical reliability when recordings are crowdsourced, labels are heterogeneous, acquisition protocols vary over time, and non-audio metadata are correlated with disease status. This paper presents a protocol-aware validation study for COVID-19 respiratory-audio modeling. Coswara is used for multimodal source-domain analysis over cough, breath, and speech; COUGHVID is used only as an independent cough-only external transfer target. The validation ladder includes participant-disjoint internal testing, time-stratified participant testing, strict early-to-late temporal testing, cough-only external transfer, metadata-only negative controls, and paper-comparable cross-validation. The strongest Coswara participant-split fusion model achieved AUROC 0.897, but performance decreased under stricter validation: time-stratified validation reached AUROC 0.849 and strict early-to-late temporal validation reached AUROC 0.698. In model-matched cough-only comparisons using frozen ComParE+IS10 handcrafted features, Coswara internal AUROC was 0.849-0.868, whereas COUGHVID transfer AUROC was 0.523-0.543. Metadata-only models achieved AUROC 0.964, with symptoms-only and demographic/protocol-only configurations reaching AUROC 0.932 and 0.914. Inverse-propensity sensitivity reduced quality-weighted audio AUROC from 0.879 unweighted to 0.780-0.807 after measured adjustment. These results show that respiratory-audio studies should report validation target, modality scope, external portability, calibration context, and metadata shortcut risk before interpreting internal performance as screening readiness.
\end{abstract}

\begin{IEEEkeywords}
COVID-19, respiratory audio, cough analysis, multimodal fusion, temporal validation, dataset shift, confounding, biomedical machine learning.
\end{IEEEkeywords}

\section{Introduction}
\IEEEPARstart{R}{espiratory-audio} machine learning has been proposed as a low-cost and non-invasive complement to conventional COVID-19 testing. The central appeal is operational: cough, breathing, and speech can be collected remotely, repeated frequently, and analyzed without specialized clinical hardware. Public datasets such as Coswara, COUGHVID, and COVID-19 Sounds have made this problem accessible to biomedical informatics researchers by coupling respiratory recordings with COVID-19 labels and participant metadata \cite{bhattacharya2023coswara,orlandic2021coughvid,xia2021covid}.

The main scientific difficulty is that respiratory-audio COVID-19 screening is not only a pattern-recognition task. Crowdsourced recordings differ by device, environment, country, collection month, symptom status, label source, and participant behavior. These factors can create dataset-specific shortcuts that are stable under random or weakly controlled splits but unstable under chronological or external validation. A model may therefore achieve high internal AUROC while learning acquisition-period structure, symptom proxies, or cohort artifacts rather than a durable acoustic correlate of infection.

This problem is especially important for biomedical AI papers because an internal classification number is easy to report and difficult to interpret. Participant-disjoint validation reduces direct identity leakage, but it does not by itself answer whether performance survives calendar shift, label-prevalence shift, collection-protocol shift, or metadata confounding. For respiratory-audio screening, a credible validation report must ask four questions together: does the model separate participants in the source dataset, does it remain stable across time, does a cough branch transfer to an independent cough corpus, and can non-audio variables alone explain much of the apparent signal?

We therefore treat validation design as part of the biomedical model, not as an afterthought. A classifier that performs well only in the source collection can still be scientifically useful, but its claim level is source-domain discrimination rather than deployable screening. Conversely, a lower score under temporal or external validation is not merely a negative result; it identifies the validation boundary where a clinical informatics claim stops being supported.

We address these questions through a protocol-aware validation study. Coswara supports multimodal source-domain and temporal analyses because cough, breath, and speech recordings are available for the same source collection. COUGHVID is used as a separate cough-only external transfer target; it is not evidence for multimodal fusion because matched breath and speech modalities are absent. The study deliberately treats negative and degraded results as reliability evidence rather than as failures to be hidden.

The contributions are as follows.
\begin{enumerate}
\item We define a validation ladder for COVID-19 respiratory-audio screening that combines participant-disjoint internal validation, time-stratified validation, strict early-to-late temporal validation, external cough transfer, and metadata-only confounding controls.
\item We evaluate a broad but controlled model bank spanning MFCC-style baselines, compact CNN branches, WavLM self-supervised representations, OpenSMILE ComParE+IS10 features, gradient-boosted and support-vector classifiers, feature-level fusion, and late fusion.
\item We show that Coswara internal multimodal performance can be strong (AUROC 0.897), while modality-matched cough-only transfer drops from Coswara AUROC 0.849-0.868 to COUGHVID AUROC 0.523-0.543.
\item We quantify metadata shortcut risk, showing that non-audio metadata alone reaches AUROC 0.964, with symptoms-only and demographic/protocol-only baselines also highly predictive.
\end{enumerate}

\section{Related Work}
\subsection{Respiratory-Audio Dataset Resources}
COVID-19 respiratory-audio modeling has been shaped by several public datasets. Coswara provides crowdsourced cough, breath, speech, symptom, and demographic/protocol information for remote SARS-CoV-2 screening research \cite{bhattacharya2023coswara}. COUGHVID provides a large cough-only collection designed for cough analysis and external benchmarking \cite{orlandic2021coughvid}. COVID-19 Sounds broadened the field with large-scale breathing, cough, and voice collection for digital respiratory screening \cite{xia2021covid}. These resources enabled rapid modeling work, but their collection processes also make validation design central to interpretation.

\subsection{Architecture-Forward COVID Audio Modeling}
Many COVID-audio papers have emphasized model architecture or representation learning. Smartphone cough classification studies reported high internal performance with handcrafted and learned representations \cite{pahar2021global}. Deep transfer learning over cough, breath, and speech motivated multimodal respiratory-audio classification \cite{pahar2022transfer}. Ensemble and multi-criteria decision methods were explored for cough-based detection \cite{chowdhury2022mcdm}. In JBHI, the Hierarchical Spectrogram Transformer (HST) demonstrated an architecture-forward strategy for respiratory-sound COVID-19 detection \cite{aytekin2024hst}. Our study is complementary: the primary contribution is not a new neural architecture, but a reliability-oriented validation framework that stresses temporal, external, and confounding-aware behavior.

\subsection{Cross-Dataset and Temporal Validation}
Cross-dataset evaluation is a stricter test than within-dataset cross-validation because collection procedures, labels, and prevalence can change simultaneously. Recent cross-dataset work with deep neural decision trees and forests reported substantially different performance across cough datasets, including limited Coswara--COUGHVID transfer in some directions \cite{islam2026dndf}. Temporal drift has also been studied directly for COVID cough audio; drift-adaptive work reported degradation from development to postdevelopment periods and motivated monitoring and adaptation strategies \cite{ganitidis2025drift}. These studies support treating chronological and external validation as core evidence rather than optional sensitivity analyses.

\subsection{Confounding and Shortcut Learning}
Confounding is a persistent risk in medical audio. Realistic evaluations of COVID-19 Sounds emphasized that audio-based digital testing must be interpreted under deployment-like constraints \cite{han2022sounds}. A recent confounding critique found that audio classifiers did not clearly improve over simple symptom checkers under realistic evaluation, highlighting the need to compare audio models against symptom and metadata baselines \cite{coppock2024audio}. In this paper, metadata-only models are not proposed as screening tools; they are negative controls that quantify how much label information is available without audio.

\begin{table*}[t]
\caption{Evidence-type positioning relative to representative COVID respiratory-audio literature. The comparison is protocol-oriented rather than a raw leaderboard because datasets, modalities, labels, and validation targets differ across studies.}
\label{tab:jbhi_literature}
\centering
\scriptsize
\begin{tabularx}{\textwidth}{p{0.15\textwidth} p{0.17\textwidth} p{0.22\textwidth} Y}
\toprule
Evidence type and study & Modality/data & Validation emphasis & Relevance to this work \\
\midrule
Dataset resource: Coswara \cite{bhattacharya2023coswara} & Cough, breath, speech, metadata & Dataset construction & Source dataset for multimodal, temporal, and metadata-control analyses. \\
Dataset resource: COUGHVID \cite{orlandic2021coughvid} & Cough & Dataset construction & Independent cough-only target for external transfer. \\
Dataset resource: COVID-19 Sounds \cite{xia2021covid} & Cough, breath, voice & Large-scale data collection & Establishes the broader digital respiratory-screening setting. \\
Architecture-forward: Pahar et al. \cite{pahar2021global,pahar2022transfer} & Smartphone cough, breath, speech & Source-domain modeling & Motivates respiratory-audio and multimodal modeling, but not by itself temporal or external reliability. \\
Architecture-forward: HST \cite{aytekin2024hst} & Respiratory sounds & Transformer modeling & Advanced JBHI comparator; our emphasis is protocol-forward validation rather than new architecture. \\
Cross-dataset: DNDT/DNDF \cite{islam2026dndf} & Multiple cough datasets & Cross-dataset evaluation & Context for weak Coswara--COUGHVID transfer and protocol-dependent comparisons. \\
Temporal drift: Ganitidis et al. \cite{ganitidis2025drift} & COVID cough audio & Chronological degradation and adaptation & Supports strict early-to-late testing as a main validation endpoint. \\
Confounding: Coppock et al. \cite{coppock2024audio} & Vocal audio, symptoms & Symptom and shortcut controls & Motivates metadata-only, symptom-only, and demographic/protocol-only negative controls. \\
\bottomrule
\end{tabularx}
\end{table*}

\section{Materials}
\subsection{Datasets and Modality Scope}
Coswara is the source dataset for the multimodal portion of this study. It contains cough, breathing, speech-like recordings, symptom variables, and demographic/protocol metadata collected for remote SARS-CoV-2 screening research \cite{bhattacharya2023coswara}. Endpoint sample sizes vary by modality availability, feature completeness, label availability, and validation protocol; therefore, each results table reports the corresponding evaluation sample size.

COUGHVID is the external dataset. It is used only for cough-domain transfer because it does not provide matched breath and speech recordings for the same multimodal endpoint \cite{orlandic2021coughvid}. Consequently, the Coswara fusion results and the COUGHVID result answer different questions: source-domain multimodal discrimination versus independent cough-domain transfer.

\subsection{Labels and Metadata}
The classification target is the dataset-provided binary COVID-19 status available for the analytic endpoint. Metadata variables are not used as evidence of an audio biomarker. Instead, they are reserved for audit analyses, including symptoms-only, demographic/protocol-only, and full safe metadata baselines. The metadata experiments test whether non-audio variables can predict labels strongly enough to threaten interpretation of audio-only and multimodal results.

\begin{table*}[t]
\caption{Dataset roles and validation protocols. Coswara supports multimodal source-domain and temporal analyses; COUGHVID is restricted to cough-only external transfer. Endpoint-specific sample sizes are reported with the corresponding results because feature completeness and validation protocol determine the final analytic set.}
\label{tab:dataset_protocols}
\centering
\scriptsize
\begin{tabularx}{\textwidth}{p{0.12\textwidth} p{0.19\textwidth} p{0.18\textwidth} p{0.17\textwidth} Y}
\toprule
Dataset & Modalities used here & Validation role & Unit of separation & Interpretation \\
\midrule
Coswara & Cough, breath, speech, metadata & Participant-disjoint internal split & Participant & Tests source-domain discrimination without participant overlap. \\
Coswara & Cough, breath, speech, metadata & Time-stratified participant split & Participant and collection-time strata & Tests internal performance after preserving calendar structure. \\
Coswara & Cough, breath, speech, metadata & Strict early-to-late temporal split & Earlier-to-later recordings & Tests chronological generalization under pandemic-period shift. \\
Coswara & Metadata only & Confounding audit & Same endpoint labels, no audio & Tests whether symptoms, demographic variables, or protocol variables can predict labels without audio. \\
COUGHVID & Cough only & External transfer & Independent dataset & Tests whether a Coswara-trained cough branch transfers to a different cough collection process. \\
\bottomrule
\end{tabularx}
\end{table*}

\section{Methods}
\subsection{Study Design}
The study was designed around validation burden rather than around a single best model. Figure~\ref{fig:pipeline} summarizes the pipeline: respiratory recordings are represented through multiple feature families, modality-specific models are trained inside source-domain protocols, fusion is applied only where matched modalities are available, and final claims are filtered through temporal, external, and metadata-control checks. Figure~\ref{fig:protocol_ladder} shows the validation ladder used to interpret model reliability.

\begin{figure*}[t]
\centering
\ownedfigurebox{End-to-end validation pipeline}{../common_figures/fig_pipeline.pdf}{}
\caption{Protocol-aware validation pipeline. Coswara cough, breath, and speech recordings support source-domain unimodal and multimodal modeling. Candidate audio representations and classifiers are evaluated under participant-disjoint, time-stratified, and strict early-to-late temporal protocols. COUGHVID enters only as an independent cough-only external transfer target. Metadata-only branches are negative controls for shortcut learning rather than candidate clinical models.}
\label{fig:pipeline}
\end{figure*}

\begin{figure}[t]
\centering
\ownedfigurebox{Evaluation protocol ladder}{../common_figures/fig_evaluation_ladder.pdf}{}
\caption{Evaluation ladder used for claim control. Each level increases the reliability burden: participant-disjoint internal testing checks source-domain separation, time-stratified testing accounts for calendar structure, strict temporal testing evaluates early-to-late generalization, COUGHVID transfer tests independent cough-domain robustness, and metadata-only controls test whether non-audio shortcuts can explain label predictability.}
\label{fig:protocol_ladder}
\end{figure}

\subsection{Audio Representations and Classifiers}
The model bank included conventional MFCC-style acoustic baselines, compact CNN branches over spectrogram-style inputs, WavLM self-supervised representations \cite{chen2022wavlm}, particle-swarm feature search, and OpenSMILE ComParE+IS10 acoustic descriptors \cite{eyben2010opensmile,schuller2010is10,schuller2016compare}. The final selected validation rows use ComParE+IS10 features with source-domain training-workflow feature selection. The top-800 feature strategy was frozen before COUGHVID transfer, and COUGHVID labels were not used for feature selection. ComParE+IS10 was used for the external transfer endpoint because it is a standardized paralinguistic representation widely aligned with handcrafted-feature respiratory-audio baselines and can be applied to the target corpus without target-domain tuning. For strict chronological validation, this is a frozen source-domain representation rather than a fully prospective feature-selection protocol re-estimated only inside the early-period cohort. Candidate classifiers included logistic regression, support-vector machines, random forests, extra-trees models, LightGBM \cite{ke2017lightgbm}, XGBoost \cite{chen2016xgboost}, CatBoost \cite{prokhorenkova2018catboost}, and related ensemble learners.

The model bank is intentionally broader than the final selected row. Its purpose is to reduce dependence on a single representation choice and to establish whether heavier learned representations remove the validation problem. They did not: the final claims are driven by protocol behavior rather than by a single architecture's headline score.

The deep branches should be interpreted carefully. The practical CNN-BiGRU and WavLM rows are single-stream source-domain checks, not completed COUGHVID transfer experiments. The 0.897 AUROC row is a validation-selected Coswara multimodal fusion endpoint, whereas the WavLM and CNN rows are isolated modality branches. This difference explains why the fusion endpoint can exceed the single-stream deep rows without implying that deep models were externally validated.

\subsection{Fusion Definitions}
Fusion was evaluated only within Coswara, where cough, breath, and speech recordings can be aligned to the source-domain endpoint. Feature-level fusion concatenates modality-specific feature vectors before classifier training. Late fusion combines modality-branch predictions after branch-level model fitting. For a participant or recording $i$ with modality set $\mathcal{M}$, the uniform late-fusion probability is
\begin{equation}
\hat{p}_{i}^{\mathrm{uniform}} =
\frac{1}{|\mathcal{M}|}\sum_{m \in \mathcal{M}}\hat{p}_{im},
\end{equation}
where $\hat{p}_{im}$ is the predicted probability from modality branch $m$. Stacked logistic fusion fits a logistic meta-model on validation-available branch predictions and applies the learned combiner to held-out predictions. COUGHVID is excluded from multimodal fusion because only cough recordings are available in the external target.

\subsection{Model Selection and Thresholding}
Model and fusion selection was performed inside the relevant validation protocol. Threshold-independent metrics were used for ranking whenever possible, and operating-point metrics were computed at validation-selected thresholds or source-selected thresholds for external transfer. This distinction matters for COUGHVID: a threshold selected on the Coswara source distribution may be poorly calibrated under the external target prevalence and label distribution.

\subsection{Evaluation Metrics}
AUROC and AUPRC are the primary threshold-independent metrics. Balanced accuracy, F1 score, sensitivity, and specificity describe the selected operating point. Balanced accuracy is
\begin{equation}
\mathrm{BA}=\frac{1}{2}\left(\frac{\mathrm{TP}}{\mathrm{TP}+\mathrm{FN}}+
\frac{\mathrm{TN}}{\mathrm{TN}+\mathrm{FP}}\right).
\end{equation}
For external transfer, AUPRC is interpreted relative to target prevalence because positive labels are rare in the external cough endpoint. The primary ladder rows are reported as point estimates. Paper-comparable cross-validation rows include fold-level variability, but structural confidence intervals were not computed for every final selected participant, temporal, and external-transfer endpoint. The ladder is therefore used to support large validation-pattern changes, not small pairwise model rankings.

\section{Evaluation Protocols}
\subsection{Participant-Disjoint Internal Validation}
The participant-disjoint internal split tests whether models discriminate COVID-19 status within the source dataset without participant overlap. This endpoint supports multimodal Coswara fusion and is the closest setting to many internal benchmark reports.

\subsection{Time-Stratified Participant Validation}
The time-stratified participant split preserves participant separation while accounting for collection-time structure. It asks whether source-domain performance remains strong when calendar imbalance is reduced rather than ignored.

\subsection{Strict Early-to-Late Temporal Validation}
Strict temporal validation trains on earlier recordings and evaluates on later recordings. This is the main chronological stress test and is interpreted as evidence about temporal robustness, not as a claim of prospective clinical validation.

\subsection{COUGHVID Cough-Only External Transfer}
The external transfer endpoint evaluates a Coswara-trained cough model on COUGHVID. Because COUGHVID is cough-only in this study, this protocol cannot validate multimodal fusion. It tests a narrower but important question: whether a cough branch trained in one crowdsourced collection process transfers to another.

\subsection{Metadata-Only Confounding Audit}
The metadata audit trains models without audio. Three configurations are reported: full safe metadata, symptoms only, and demographic/protocol only. A strict temporal month-ablation is also reported to assess whether explicit calendar variables behave as shortcuts under early-to-late evaluation.

\begin{table*}[t]
\caption{Claim-control interpretation used throughout the study. Each protocol answers a different biomedical question; the admissible claim is deliberately narrower than the best internal metric when later validation tiers fail.}
\label{tab:claim_control}
\centering
\scriptsize
\begin{tabularx}{\textwidth}{p{0.19\textwidth} p{0.26\textwidth} p{0.19\textwidth} Y}
\toprule
Validation evidence & Question answered & Main observed result & Claim boundary \\
\midrule
Participant-disjoint Coswara fusion & Can source-domain cough, breath, or speech models discriminate labels without participant overlap? & Best selected fusion AUROC 0.897 & Supports source-domain feasibility, not future-period or external reliability. \\
Time-aware and strict temporal testing & Does performance survive collection-time structure and early-to-late shift? & AUROC 0.849 under time stratification; 0.698 under strict temporal testing & Supports a temporal degradation claim and requires caution for deployment-like use. \\
COUGHVID cough transfer & Does a Coswara-trained cough branch transfer to an independent cough corpus? & AUROC 0.543 and AUPRC 0.040 & Supports weak cough portability; does not test multimodal external fusion. \\
Metadata-only controls & Can measured non-audio variables recover labels without audio? & Full metadata AUROC 0.964; demographic/protocol AUROC 0.914 & Supports shortcut-risk and confounding interpretation, not metadata deployment. \\
\bottomrule
\end{tabularx}
\end{table*}

\section{Results}
\subsection{Internal Fusion Is Strong Under Source-Domain Validation}
The strongest internal Coswara endpoint was a cough+speech ComParE+IS10 stacked logistic fusion model with AUROC 0.897 and AUPRC 0.863 on 314 evaluation samples (Table~\ref{tab:primary_results}). Balanced accuracy was 0.825, with sensitivity 0.833 and specificity 0.816. This result shows that source-domain multimodal respiratory audio can carry discriminative signal when participants are separated.

\subsection{Time-Aware and Strict Temporal Validation Reduce Performance}
The stronger reliability tests reduced performance. The time-stratified cough+breath+speech uniform-mean fusion row reached AUROC 0.849 and AUPRC 0.783 on 431 samples. Under strict early-to-late temporal validation, the selected breath ComParE+IS10 top-4 ensemble reached AUROC 0.698 on 411 samples. Figure~\ref{fig:temporal_degradation} summarizes this degradation. The temporal endpoint retains nontrivial discrimination, but it is materially weaker than the internal participant-split result and should not be interpreted as deployment-ready performance.

\subsection{COUGHVID External Transfer Is Near Chance and Cough-Only}
The COUGHVID endpoint is a cough-only external transfer test. The selected Coswara-to-COUGHVID cough model reached AUROC 0.543 and AUPRC 0.040 on 8331 target samples. Balanced accuracy was 0.510, F1 was 0.058, sensitivity was 0.084, and specificity was 0.936. The high specificity with very low sensitivity indicates an operating-point mismatch under external shift. The AUPRC is only slightly above the target prevalence reported in the support analysis (0.034), reinforcing that the external cough transfer result provides weak evidence of generalization.

The selected rows in Table~\ref{tab:primary_results} are claim-control endpoints rather than a fixed-model ablation; the selected modality and fusion rule change across tiers. The fair comparison for COUGHVID is cough-only on both sides. Table~\ref{tab:cough_matched_transfer} therefore compares Coswara cough-only internal rows against the corresponding COUGHVID cough-only transfer rows under the same ComParE+IS10 top-800 feature strategy. Coswara cough-only internal AUROC ranged from 0.849 to 0.868 across the matched model families, whereas COUGHVID transfer AUROC ranged from 0.523 to 0.543. Table~\ref{tab:fixed_cough_ladder} extends the matched comparison across all validation tiers while holding modality fixed to cough and model family fixed. The external degradation is therefore not only a consequence of missing breath and speech modalities; the matched cough branch itself is weakly portable.

\begin{table*}[t]
\caption{Selected audio endpoints across the validation ladder. Values are rounded to three decimals and are interpreted by protocol, not as a single pooled leaderboard. Coswara rows can involve multimodal fusion; the COUGHVID row is cough-only external transfer.}
\label{tab:primary_results}
\centering
\scriptsize
\resizebox{\textwidth}{!}{%
\begin{tabular}{p{0.23\textwidth} p{0.19\textwidth} p{0.20\textwidth} c c c c c c c}
\toprule
Protocol & Modality/model & Fusion or selection & $n$ & AUROC & AUPRC & Bal. Acc. & F1 & Sens. & Spec. \\
\midrule
Existing participant split & Cough+speech ComParE+IS10 & Stacked logistic fusion & 314 & 0.897 & 0.863 & 0.825 & 0.752 & 0.833 & 0.816 \\
Time-stratified participant split & Cough+breath+speech ComParE+IS10 & Uniform mean & 431 & 0.849 & 0.783 & 0.783 & 0.705 & 0.710 & 0.857 \\
Strict early-to-late temporal & Breath ComParE+IS10 & Top-4 validation ensemble & 411 & 0.698 & 0.896 & 0.656 & 0.751 & 0.646 & 0.667 \\
COUGHVID external transfer & Cough ComParE+IS10 LightGBM & Source validation threshold & 8331 & 0.543 & 0.040 & 0.510 & 0.058 & 0.084 & 0.936 \\
\bottomrule
\end{tabular}
}
\end{table*}

\begin{table*}[t]
\caption{Model-matched cough-only internal-to-external comparison. Internal rows use Coswara participant-disjoint cough endpoints with ComParE+IS10 top-800 features. External rows evaluate the corresponding Coswara-trained cough model on COUGHVID.}
\label{tab:cough_matched_transfer}
\centering
\scriptsize
\begin{tabular}{l ccc ccc}
\toprule
Model & \multicolumn{3}{c}{Coswara cough-only} & \multicolumn{3}{c}{COUGHVID cough-only} \\
\cmidrule(lr){2-4}\cmidrule(lr){5-7}
 & AUROC & AUPRC & Bal. Acc. & AUROC & AUPRC & Bal. Acc. \\
\midrule
LightGBM SMOTE f80 & 0.853 & 0.797 & 0.782 & 0.543 & 0.040 & 0.510 \\
SVC RBF f60 & 0.868 & 0.812 & 0.795 & 0.523 & 0.037 & 0.500 \\
CatBoost SMOTE f80 & 0.849 & 0.788 & 0.776 & 0.531 & 0.040 & 0.506 \\
XGBoost SMOTE f80 & 0.850 & 0.780 & 0.776 & 0.532 & 0.039 & 0.505 \\
\bottomrule
\end{tabular}
\end{table*}

\begin{table*}[t]
\caption{Fixed cough-only model-family ladder. Each row keeps the modality fixed to cough and follows the same model family across validation tiers. Models are retrained within the corresponding source-domain protocol; values are AUROC.}
\label{tab:fixed_cough_ladder}
\centering
\scriptsize
\begin{tabular}{lcccc}
\toprule
Model & Existing participant split & Time-stratified split & Strict temporal split & COUGHVID external \\
\midrule
LightGBM SMOTE & 0.853 & 0.830 & 0.604 & 0.543 \\
SVC RBF & 0.868 & 0.802 & 0.561 & 0.523 \\
CatBoost SMOTE & 0.849 & 0.826 & 0.614 & 0.531 \\
XGBoost SMOTE & 0.850 & 0.828 & 0.599 & 0.532 \\
\bottomrule
\end{tabular}
\end{table*}

\begin{figure}[t]
\centering
\ownedfigurebox{Validation degradation}{../common_figures/fig_validation_degradation.pdf}{}
\caption{AUROC degradation across validation settings. The plotted endpoints are the selected rows from Table~\ref{tab:primary_results}: Coswara participant-disjoint fusion, Coswara time-stratified fusion, Coswara strict temporal validation, and COUGHVID cough-only external transfer. Because these selected endpoints can differ in modality and fusion rule, Table~\ref{tab:fixed_cough_ladder} reports the modality- and model-family-matched cough-only ladder. Together, the rows show why an internal respiratory-audio score is not sufficient evidence of external reliability.}
\label{fig:temporal_degradation}
\end{figure}

\subsection{Metadata-Only Baselines Reveal Strong Confounding Risk}
Metadata-only models were highly predictive in internal validation (Table~\ref{tab:metadata}). Full safe metadata reached AUROC 0.964, AUPRC 0.928, and balanced accuracy 0.890 on 2862 samples. Symptoms-only metadata reached AUROC 0.932, and demographic/protocol-only metadata reached AUROC 0.914. These values are high enough to challenge any interpretation that the audio models isolate a stable disease-specific acoustic signature.

The temporal metadata ablation provides an additional warning. In strict temporal validation, full metadata reached AUROC 0.531, while removing recording month increased AUROC to 0.779. This increase of 0.247 suggests that explicit calendar variables can encode source-period associations that do not transport cleanly to later data. The effect should be interpreted as shortcut-risk evidence, not as a causal estimate of calendar time.

\subsection{IPW Sensitivity Qualifies the Audio Signal}
Inverse-propensity weighting was used as a measured-confounding sensitivity analysis for the quality-weighted audio fusion prediction set (Table~\ref{tab:ipw}). The unweighted row reached AUROC 0.879 and AUPRC 0.832 on 318 held-out samples. With a strict cap of 2, AUROC remained 0.807 and AUPRC 0.624, with effective sample size 238.7. With cap 10 and 99th-percentile truncation, AUROC was 0.780 and AUPRC 0.537, while the maximum absolute residual SMD decreased to 0.382 and effective sample size decreased to 130.4. This pattern supports a qualified conclusion: audio signal remains above chance after measured adjustment, but the drop from the unweighted row indicates that observational selection and measured imbalance materially inflate naive source-domain performance.

\begin{table}[t]
\caption{Inverse-propensity weighting sensitivity analysis for quality-weighted audio fusion. ESS denotes effective sample size; SMD denotes standardized mean difference.}
\label{tab:ipw}
\centering
\footnotesize
\resizebox{\columnwidth}{!}{%
\begin{tabular}{lrrrrrr}
\toprule
Weighting setting & AUROC & AUPRC & Bal. Acc. & ESS & Mean abs. SMD & Max abs. SMD \\
\midrule
Unweighted & 0.879 & 0.832 & 0.804 & 318.0 & 0.224 & 1.384 \\
IPW cap 2, q=0.95 & 0.807 & 0.624 & 0.742 & 238.7 & 0.149 & 0.724 \\
IPW cap 5, q=0.95 & 0.793 & 0.573 & 0.731 & 173.2 & 0.129 & 0.523 \\
IPW cap 10, q=0.99 & 0.780 & 0.537 & 0.721 & 130.4 & 0.118 & 0.382 \\
\bottomrule
\end{tabular}}
\end{table}

\begin{table}[t]
\caption{Metadata-only confounding evidence. Metadata rows are negative controls, not proposed clinical models. The strict temporal month-ablation shows that calendar variables can behave differently under chronological shift than under internal validation.}
\label{tab:metadata}
\centering
\footnotesize
\resizebox{\columnwidth}{!}{%
\begin{tabular}{l c c c c}
\toprule
Metadata setting & $n$ & AUROC & AUPRC & Bal. Acc. \\
\midrule
Full safe metadata & 2862 & 0.964 & 0.928 & 0.890 \\
Symptoms only & 2862 & 0.932 & 0.898 & 0.912 \\
Demographic/protocol only & 2862 & 0.914 & 0.737 & 0.827 \\
Strict temporal full metadata & 3816 & 0.531 & 0.837 & 0.492 \\
Strict temporal, month removed & 3816 & 0.779 & 0.935 & 0.723 \\
\bottomrule
\end{tabular}
}
\end{table}

\begin{figure}[t]
\centering
\ownedfigurebox{Metadata confounding evidence}{../common_figures/fig_metadata_confounding.pdf}{}
\caption{Metadata-confounding audit. Internal metadata-only AUROC values are high even without audio, while the strict temporal month-ablation indicates that collection-time variables can act as unstable shortcuts. The figure supports interpreting metadata models as reliability controls rather than diagnostic systems.}
\label{fig:confounding}
\end{figure}

\subsection{Paper-Comparable Cross-Validation Is Useful but Less Decisive}
Paper-comparable 10-fold cross-validation was run to connect the study to common reporting practice (Table~\ref{tab:cv}). The best cough row was an RBF-kernel support-vector classifier with AUROC $0.819 \pm 0.029$ and AUPRC $0.728 \pm 0.036$ on 4094 samples. The best speech row was LightGBM with AUROC $0.806 \pm 0.014$ and AUPRC $0.705 \pm 0.021$ on 9983 samples. These numbers are not the main reliability evidence because recording-level cross-validation is less stringent than strict temporal or external transfer evaluation.

\begin{table}[t]
\caption{Paper-comparable 10-fold cross-validation. These rows support comparison with common within-dataset reporting practice, but they are not used to claim external reliability.}
\label{tab:cv}
\centering
\footnotesize
\resizebox{\columnwidth}{!}{%
\begin{tabular}{l l c c c c c}
\toprule
Modality & Best model & $n$ & AUROC & AUPRC & Bal. Acc. & F1 \\
\midrule
Cough & SVC RBF & 4094 & $0.819 \pm 0.029$ & $0.728 \pm 0.036$ & 0.737 & 0.645 \\
Speech & LightGBM & 9983 & $0.806 \pm 0.014$ & $0.705 \pm 0.021$ & 0.732 & 0.635 \\
\bottomrule
\end{tabular}
}
\end{table}

\subsection{Heavier Branches Did Not Remove the Validation Problem}
Representative non-final rows are shown in Table~\ref{tab:negative}. CNN/spectrogram branches and WavLM branches produced plausible internal single-modality results, but they did not supersede the final ComParE+IS10 validation-selected endpoints. This is important for interpretation: the observed temporal and external fragility is not merely the result of choosing an unusually weak representation. It also prevents a misleading comparison: the WavLM and CNN rows were not multimodal fusion rows and were not completed as COUGHVID transfer rows.

\begin{table}[t]
\caption{Representative non-final model branches. Architecture or representation complexity alone did not remove the need for stricter validation.}
\label{tab:negative}
\centering
\footnotesize
\resizebox{\columnwidth}{!}{%
\begin{tabular}{l l c c c c}
\toprule
Branch & Representative row & $n$ & AUROC & AUPRC & Bal. Acc. \\
\midrule
CNN & CNN-BiGRU cough & 636 & 0.730 & 0.556 & 0.686 \\
CNN & CNN-BiGRU breath & 636 & 0.710 & 0.590 & 0.649 \\
CNN & CNN-BiGRU speech & 1590 & 0.676 & 0.490 & 0.638 \\
WavLM & Shallow breath & 296 & 0.790 & 0.654 & 0.731 \\
WavLM & Shallow cough & 310 & 0.790 & 0.731 & 0.713 \\
WavLM & Heavy cough & 312 & 0.789 & 0.696 & 0.730 \\
\bottomrule
\end{tabular}
}
\end{table}

\begin{figure}[t]
\centering
\ownedfigurebox{Multimodal late fusion}{../common_figures/fig_multimodal_fusion.pdf}{}
\caption{Modality scope of the evaluated endpoints. Multimodal fusion is restricted to Coswara because cough, breath, and speech are available within the source-domain protocol. COUGHVID contributes only cough recordings to this study, so it evaluates cough-domain external transfer rather than multimodal generalization.}
\label{fig:fusion}
\end{figure}

\section{Discussion}
The main finding is that COVID-19 respiratory-audio models can look substantially stronger under internal source-domain validation than under deployment-like stress tests. The best Coswara internal fusion row reached AUROC 0.897, but stricter validation reduced the claim: time-stratified validation reached AUROC 0.849 and strict early-to-late temporal validation reached AUROC 0.698. For external testing, the correct modality-matched comparison is Coswara cough-only internal AUROC 0.849-0.868 versus COUGHVID cough-only transfer AUROC 0.523-0.543. This pattern aligns with prior temporal-drift and cross-dataset evidence that respiratory-audio models are sensitive to collection conditions and validation design \cite{ganitidis2025drift,islam2026dndf}.

The metadata results are equally important. Full safe metadata alone reached AUROC 0.964, symptoms-only metadata reached AUROC 0.932, and demographic/protocol-only metadata reached AUROC 0.914. Symptoms may be clinically related to infection status, but a model that exploits symptom structure is not necessarily detecting acoustic pathophysiology. The demographic/protocol result is more concerning because it shows that label information is also available from non-symptom variables. Support analyses identified recording year, recording month, country, age, and recording duration among leading demographic/protocol attribution drivers, consistent with the interpretation that cohort and acquisition structure are entangled with labels.

The external transfer result should be interpreted narrowly and conservatively. It does not test multimodal fusion, because the external target is cough-only. It also does not prove that all cough models must fail under all external datasets. It does show that a strong Coswara-selected cough representation did not transport well to COUGHVID under the evaluated protocol. The low sensitivity at the source-selected operating point suggests that calibration and thresholding are not portable across the source and target distributions.

For JBHI-style biomedical informatics reporting, these findings argue for protocol-aware evidence packages. Internal AUROC remains useful, especially when participants are disjoint, but it should be accompanied by chronological testing, external testing when modality-compatible data exist, metadata-only negative controls, and clear modality scoping. Without those controls, respiratory-audio studies risk reporting source-dataset discrimination as if it were deployable screening reliability.

The practical reporting recommendation is direct. A COVID-19 respiratory-audio paper should identify the unit of separation, the time structure of the split, whether external data match the source modalities, how thresholds are selected, whether AUPRC is interpreted against prevalence, and whether symptoms or protocol metadata can predict labels without audio. These requirements are not cosmetic reporting details; they determine whether a reported model is a source-domain classifier, a future-period classifier, an externally portable classifier, or only an exploratory benchmark.

\section{Limitations and Future Validation}
The limitations define the supported claim level and the next validation requirements. First, COUGHVID is cough-only in this study. It provides an independent test of Coswara-to-COUGHVID cough transfer, but it cannot validate Coswara breath, speech, or multimodal fusion. A stronger next study would require an external cohort with matched cough, breath, and speech prompts, collected under a known protocol and evaluated without source-dataset retuning.

Second, the datasets are retrospective and crowdsourced, with heterogeneous labels, devices, environments, prompts, and recruitment processes. Participant-disjoint and temporal validation reduce some over-optimism, but they do not replace a prospective clinical screening cohort with prespecified thresholds, calibration monitoring, and clinically adjudicated endpoints. The temporal analysis is therefore a retrospective stress test, not a substitute for prospective validation.

Third, metadata and inverse-propensity analyses quantify measured shortcut risk rather than proving a complete causal pathway. Strong symptoms-only and demographic/protocol-only baselines show that non-audio information is available, but they do not identify every latent feature used by every audio model. The IPW sensitivity rows should be read as robustness checks under measured adjustment: cap-10 weighting improves residual measured balance but reduces effective sample size to 130.4, so the retained AUROC 0.780 is supportive but not definitive. Fourth, the primary ladder does not yet include DeLong, bootstrap, or paired-resampling confidence intervals for every selected endpoint. The large internal-to-temporal and internal-to-external changes motivate reliability concern, but individual small model differences should be treated as descriptive until interval estimates are added. Fifth, the model bank is broad but not exhaustive; additional transformer, adaptation, segmentation, saliency, or calibration strategies may improve source-domain performance. The biomedical question remains whether such improvements survive temporal, external, and confounding-aware validation. Finally, thresholded metrics were selected within validation protocols and should not be interpreted as clinical operating points without prospective threshold and calibration studies.

\section{Conclusion}
We presented a protocol-aware validation study for COVID-19 respiratory-audio screening under temporal, external, and metadata-confounding stress. Coswara supported multimodal source-domain and temporal analyses, while COUGHVID provided a cough-only external transfer target. The best internal fusion result was strong, but performance degraded under stricter temporal validation and approached chance under external cough transfer. Metadata-only baselines were highly predictive, reinforcing the risk of shortcut learning. The practical conclusion is that COVID-19 respiratory-audio papers should report modality-aware temporal, external, and confounding-aware validation before making deployment-oriented claims.

\section*{Acknowledgment}
The authors acknowledge the maintainers and participants of the public respiratory-audio datasets used in this validation study.

\bibliographystyle{IEEEtran}
\bibliography{references}

\end{document}
